\documentclass[dvipdfmx, uplatex]{jsarticle}
\usepackage[dvipdfmx]{graphicx}
\usepackage{amsmath}
\usepackage{amssymb}
\usepackage{amsfonts}
\usepackage{amsthm}
\usepackage{mathrsfs}
\usepackage{bm}
\usepackage{braket}
\usepackage{xr}
\externaldocument{exercise_double_well_resurgence}

\usepackage{silence}
\WarningFilter*{hyperref}{Token not allowed in a PDF string}
\usepackage[dvipdfmx]{hyperref}
\usepackage{pxjahyper}
\allowdisplaybreaks[4]

\usepackage[left=25truemm,top=20truemm,bottom=20truemm,right=25truemm]{geometry}

\newcommand{\hb}{\hbar}
\newcommand{\Bor}{\mathcal{B}}
\newcommand{\cE}{\mathcal{E}}

\theoremstyle{definition}
\newtheorem*{sol}{解答}

\title{解答：二重井戸の摂動級数とインスタントン背景での Borel 和\\
\large ―― \texttt{exercise\_double\_well\_resurgence.tex} の各問の解答 ――}
\author{}
\date{}

\begin{document}
\maketitle
\vspace{-6zh}

記号・式番号は問題ファイル \texttt{exercise\_double\_well\_resurgence.tex} のものを踏襲する
（\texttt{xr} により本文中の \eqref{eq:BW} 等は問題ファイルの式番号を指す）。
\textbf{各問の数値解（採点用）}は太字で示す。

\tableofcontents

%======================================================================
\section{問題 1：古典解}
%======================================================================

\begin{sol}
\textbf{(1) 真空解。}$V'(x)=\tfrac12 x(x^2-1)$ ゆえ $V'(\pm1)=0$，定数 $x\equiv\pm1$ は $\ddot x=0=V'$ を満たす。
$V(\pm1)=0,\ \dot x=0$ ゆえ作用 $S_E=0$。2 つの縮退真空。

\textbf{(2) インスタントン。}$\tfrac12\dot x^2=V=\tfrac18(1-x^2)^2$ より $\dot x=\tfrac12(1-x^2)$。
$\int\frac{dx}{1-x^2}=\operatorname{arctanh}x=\tfrac12(\tau-\tau_0)$ ゆえ
$\boxed{x_{\mathrm{cl}}(\tau)=\tanh\tfrac{\tau-\tau_0}{2}}$。
$\tau\to+\infty$ で $\to+1$，$\tau\to-\infty$ で $\to-1$（境界条件を極限の意味で満たす弱解）。
$\tanh z=1-2e^{-2z}+\cdots$ より $x_{\mathrm{cl}}\to1-2e^{-(\tau-\tau_0)}$（$\omega=1$）。反インスタントンは $-\tanh\tfrac{\tau-\tau_0}{2}$。

\textbf{(3) 作用。}
\[
    S_0=\int_{-1}^{1}\tfrac12(1-x^2)\,dx=\tfrac12\Bigl[x-\tfrac{x^3}{3}\Bigr]_{-1}^{1}=\tfrac12\Bigl(\tfrac23+\tfrac23\Bigr)
    =\boxed{\dfrac23}.
\]

\textbf{(4) N-インスタントン弱解。}よく離れた $N$ 個の(反)インスタントンの作用は，相互作用
（$\propto e^{-|\tau_{i+1}-\tau_i|}$）が分離 $\to\infty$ で消えるため，個々の和 $\to\boxed{N S_0}$。
両端で同じ真空へ戻る配位（基底状態の振幅に効く）は I と $\bar{\mathrm I}$ が偶数個 $2N$ 個並んだもので，
作用は $2NS_0$。
\end{sol}

%======================================================================
\section{問題 2：無摂動基底状態}
%======================================================================

\begin{sol}
\textbf{(1)} $[a,a^\dagger]=\tfrac12([u,-\partial_u]+[\partial_u,u])=\tfrac12(1+1)=1$。
$a^\dagger a=\tfrac12(u^2-\partial_u^2-1)=H_0-\tfrac12$ ($\partial_u u=1+u\partial_u$) ゆえ $H_0=a^\dagger a+\tfrac12$。

\textbf{(2)} $a\psi_0=0$: $(u+\partial_u)\psi_0=0\Rightarrow\psi_0=Ce^{-u^2/2}$。規格化 $C^2\sqrt\pi=1$ ゆえ
$\boxed{\psi_0=\pi^{-1/4}e^{-u^2/2}}$。$H_0\psi_0=\tfrac12\psi_0$ ゆえ
$\boxed{\cE_0=\tfrac12}$，次元を戻すと $\boxed{a_1=\tfrac12}$（零点・1-loop energy $\tfrac\hb2$）。

\textbf{(3)} $\psi_0$ は $x=1$ の単一井戸に局在し，トンネルを含まない。
\end{sol}

%======================================================================
\section{問題 3：Bender--Wu 漸化式}
%======================================================================

\begin{sol}
\textbf{(1)} $\partial_u^2(e^{-u^2/2}P)=e^{-u^2/2}(P''-2uP'+(u^2-1)P)$ より
$(-\tfrac12\partial_u^2+\tfrac12u^2)(e^{-u^2/2}P)=e^{-u^2/2}(-\tfrac12P''+uP'+\tfrac12P)$。
非調和項を加え $e^{-u^2/2}$ で割って \eqref{eq:Peq}。$P=1$ で左辺 $\tfrac12=$ 右辺 $\cE$，$\cE=\tfrac12$。

\textbf{(2)} $\epsilon^k$ の比較で \eqref{eq:BW}：
$-\tfrac12P_k''+uP_k'=-\tfrac12u^3P_{k-1}-\tfrac18u^4P_{k-2}+\sum_{j=1}^k\cE_jP_{k-j}$。
\end{sol}

%======================================================================
\section{問題 4：低次係数}
%======================================================================

\begin{sol}
左辺は $u^m\mapsto m\,u^m-\tfrac12m(m-1)u^{m-2}$。$P_k=\sum_m c_mu^m$ の $u^m$ 係数は
$mc_m-\tfrac12(m+2)(m+1)c_{m+2}$。

\textbf{(1)} $u^0$ の式は $-c_2=[\text{右辺の }u^0]=\cE_k$（右辺の $u^0$ 項は $\cE_kP_0=\cE_k$）。ゆえ
$\cE_k=-[P_k\text{ の }u^2\text{ 係数}]$。

\textbf{(2) $k=1$.} 右辺 $-\tfrac12u^3+\cE_1$。$P_1=au^3+bu$ の左辺 $3au^3+(b-3a)u$。
$3a=-\tfrac12,\ b=3a$，$u^0:\cE_1=0$。$\boxed{P_1=-\tfrac16u^3-\tfrac12u,\ \cE_1=0}$。

\textbf{(3) $k=2$.} 右辺 $=-\tfrac12u^3P_1-\tfrac18u^4+\cE_2=\tfrac1{12}u^6+\tfrac18u^4+\cE_2$。
$P_2=cu^6+du^4+eu^2$ の左辺 $6cu^6+(4d-15c)u^4+(2e-6d)u^2-e$。
$6c=\tfrac1{12}\Rightarrow c=\tfrac1{72}$；$4d-15c=\tfrac18\Rightarrow d=\tfrac1{12}$；$2e-6d=0\Rightarrow e=\tfrac14$；$-e=\cE_2$。
$\boxed{P_2=\tfrac1{72}u^6+\tfrac1{12}u^4+\tfrac14u^2,\ a_2=\cE_2=-\tfrac14}$。

\textbf{(4) $k=3,4$.} $P_3$ は奇関数ゆえ $\boxed{\cE_3=0}$，
$P_3=-\tfrac1{1296}u^9-\tfrac1{144}u^7-\tfrac1{24}u^5-\tfrac18u^3-\tfrac14u$。
$k=4$：右辺 $=-\tfrac12u^3P_3-\tfrac18u^4P_2+\cE_2P_2+\cE_4$ の各冪
$r_{12}=\tfrac1{2592},r_{10}=\tfrac1{576},r_8=\tfrac1{96},r_6=\tfrac1{36},r_4=\tfrac5{48},r_2=-\tfrac1{16}$。
$mc_m-\tfrac12(m+2)(m+1)c_{m+2}=r_m$ を上から解くと
$c_{12}=\tfrac1{31104},c_{10}=\tfrac1{2592},c_8=\tfrac1{288},c_6=\tfrac1{48},c_4=\tfrac5{48},c_2=\tfrac9{32}$，
$-c_2=\cE_4$ ゆえ $\boxed{a_3=\cE_4=-\tfrac9{32}}$。（続けて $a_4=\cE_6=-\tfrac{89}{128},\ a_5=\cE_8=-\tfrac{5013}{2048}$。）
\end{sol}

%======================================================================
\section{問題 5：Gevrey-1 発散と収束半径}
%======================================================================

\begin{sol}
\textbf{(1) 構造的起源。}$\deg P_k=3k$，下降求解 $c_m=\tfrac1m[r_m+\tfrac12(m+2)(m+1)c_{m+2}]$ で
高次係数が因子 $\tfrac{(m+2)(m+1)}{2m}\sim\tfrac m2$ を掛けながら $u^2$ 係数へ降りる。
インデックス比例因子の蓄積が $\cE_k=-c_2$ の階乗的増大を生む（定量化は (2)）。

\textbf{(2) 大次数解析。}$\operatorname{Im}E_0\sim-Ce^{-2S_0/\hb}$ を分散関係
$a_n=\tfrac1\pi\int_0^\infty\frac{\operatorname{Im}E_0}{\hb^{n+1}}d\hb$ に入れ，$t=1/\hb$ で
$\int_0^\infty t^{n-1}e^{-2S_0t}dt=\Gamma(n)/(2S_0)^n$ ゆえ
\[
    \boxed{\;a_n\sim-\frac{C}{\pi}\frac{\Gamma(n)}{(2S_0)^n}=-\frac{C}{\pi}\frac{(n-1)!}{(2S_0)^n}\;}\quad(n\to\infty).
\]
$(n-1)!$ 並みの増大＝Gevrey-1。スケールは I--$\bar{\mathrm I}$ 作用 $2S_0$。

\textbf{(3) 収束半径。}$\bigl|a_{n+1}/a_n\bigr|=\tfrac{\Gamma(n+1)}{\Gamma(n)}\tfrac1{2S_0}=\tfrac{n}{2S_0}\to\infty$，
ゆえ $\boxed{R=\lim|a_n/a_{n+1}|=\lim\tfrac{2S_0}{n}=0}$。収束半径は導いた高次挙動の帰結。

\textbf{(4) チェック.}$|a_2/a_1|=\tfrac{1/4}{1/2}=0.5$, $|a_3/a_2|=\tfrac{9/32}{1/4}=1.125$。
予言 $n/(2S_0)=0.75n$ は $n=1,2$ で $0.75,1.5$。$1/n$ 補正の範囲で矛盾しない（$2S_0=4/3$ を先取り）。
\end{sol}

%======================================================================
\section{問題 6：Borel 変換}
%======================================================================

\begin{sol}
$b_n=a_n/(n-1)!\sim K(2S_0)^{-n}$ は公比 $1/(2S_0)$ の等比数列ゆえ $|\zeta|<2S_0$ で
\[
    \Bor[E_0](\zeta)=\sum_{n\ge1}b_n\zeta^{n-1}
    =\frac{K}{2S_0}\frac{1}{1-\zeta/2S_0}=\boxed{\frac{K}{2S_0-\zeta}}.
\]
最近接特異点は\textbf{正実軸上 $\zeta=2S_0$ の単純極}（真の特異点は対数分岐だが，位置と主導的増大率は単純極で
正しく captured される）。極が積分路 $(0,\infty)$ 上にあるため，横方向 Borel 和の差が
留数 $\propto e^{-2S_0/\hb}$ の虚数的曖昧さ（I--$\bar{\mathrm I}$ 非摂動効果）を生む。
\end{sol}

%======================================================================
\section{問題 7：周期的 Borel 特異性}
%======================================================================

\begin{sol}
\textbf{(1)} Borel 平面の特異点 $\zeta=A$ は非摂動因子 $e^{-A/\hb}$ に対応する。
$2N$ 個（N 対）の多重インスタントン弱解は作用 $2NS_0$，振幅 $\propto e^{-2NS_0/\hb}$ ゆえ，
$\Bor[E_0]$ は\textbf{正実軸上の点列 $\boxed{\zeta=2NS_0\ (N=1,2,3,\dots)}$ に特異性を持つ}。
これらは\textbf{等間隔（周期的），間隔 $2S_0$}。$S_0=\tfrac23$ を入れると最初の 3 点は
$\boxed{\zeta=\tfrac43,\ \tfrac83,\ 4}$。

\textbf{(2)} $N$ 番目は $a_n$ に $\propto\Gamma(n)(2NS_0)^{-n}$ を与える。$N=2$ の $N=1$ に対する相対抑制は
$\boxed{(2S_0/4S_0)^n=(1/2)^n=2^{-n}}$。ゆえ大次数は最近接 $N=1$（$2S_0$）が支配し，
多重インスタントンは指数的に抑制された subleading 補正を与える。

\textbf{(3)} 古典 N-インスタントン弱解（作用 $2NS_0$）$\leftrightarrow$ 量子 Borel 平面の周期的特異点 $\zeta=2NS_0$ が
一対一対応する。（参考：全特異点を単純極でモデル化すると $\Bor\sim\sum_N r_N/(2NS_0-\zeta)$ で，等間隔極は
digamma 型関数の極構造に対応する。）
\end{sol}

%======================================================================
\section{問題 8：低次 Borel--Pad\'e}
%======================================================================

\begin{sol}
\textbf{(1)} $b_n=a_n/(n-1)!$：$\boxed{b_1=\tfrac12,\ b_2=-\tfrac14,\ b_3=-\tfrac9{64}}$
（$c_0=b_1,c_1=b_2,c_2=b_3$）。

\textbf{(2) $[1/1]$.} \eqref{eq:padeden} の $L=M=1,\ k=1$：$c_2+q_1c_1=0$ ゆえ $q_1=-c_2/c_1=-b_3/b_2$。
極 $\zeta^\ast_{[1/1]}=-1/q_1=b_2/b_3=\dfrac{-1/4}{-9/64}=\boxed{\dfrac{16}{9}\approx1.78}$。
最近接特異点 $2S_0=\tfrac43\approx1.333$ をやや上回る。

\textbf{(3) $[2/2]$.} $\boxed{b_4=-\tfrac{89}{768},\ b_5=-\tfrac{5013}{49152}}$。
\eqref{eq:padeden} の $L=M=2,\ k=1,2$：$b_4+q_1b_3+q_2b_2=0,\ b_5+q_1b_4+q_2b_3=0$。
整理（$\times(-64),\times(-768)$）して
\[
    9q_1+16q_2=-\tfrac{89}{12},\qquad 89q_1+108q_2=-\tfrac{5013}{64}.
\]
解いて $q_2=\tfrac{8615}{86784}\approx0.09927,\ q_1\approx-1.00055$。根
$\zeta=\dfrac{-q_1\pm\sqrt{q_1^2-4q_2}}{2q_2}=\dfrac{1.00055\pm0.77719}{0.19854}$ より
最近接 $\boxed{\zeta^\ast_{[2/2]}\approx1.12}$（他方 $\approx8.95$）。

\textbf{(4)} $[1/1]$ の $1.78$ と $[2/2]$ の $1.12$ は $2S_0=\tfrac43\approx1.333$ を上下から\textbf{挟む}。
次数を上げると最近接極は $\zeta=2S_0$ へ収束（$n\le33$ の $[15/15]$ で $\zeta^\ast\approx1.3331\approx4/3$）。
低次の数項だけで非摂動スケール $2S_0$ が当たる。
\end{sol}

%======================================================================
\section{問題 9：ループを閉じる}
%======================================================================

\begin{sol}
\textbf{(1)} 最近接特異点 $\zeta^\ast$ なら $b_n\sim(\zeta^\ast)^{-n}$，$a_n=(n-1)!\,b_n\sim\Gamma(n)(\zeta^\ast)^{-n}$。
低次から当てた $\zeta^\ast\approx4/3=2S_0$ は問題 5 の大次数則（増大率 $1/(2S_0)=3/4$）を再現する。

\textbf{(2)} $\zeta^\ast=2S_0=4/3$ は前半のインスタントン作用 $S_0=2/3$ の 2 倍。周期的特異点 $2NS_0$
（問題 7）は古典 N-インスタントン弱解に対応。無摂動基底状態への単純な補正計算に過ぎない摂動級数が，
低次の数項のうちにトンネル（多重インスタントン）の非摂動スケールを内包する
――「摂動展開から非摂動効果が拾える」共鳴（resurgence）の核心。
\end{sol}

\end{document}
